\section{Filipino Orthography and System Design}\label{sec:filipino_orthography}

The orthographic properties of Filipino directly motivate how our system
handles phonetic vowel resolution, consonant mapping, and syllable repair. This
section describes these properties and explains how each shapes the
corresponding stage of the system.

\subsection{Vowel Spelling and Phonetic Resolution}

Filipino orthography uses five vowel letters (\gr{a}, \gr{e}, \gr{i},
\gr{o}, \gr{u}), each with a stable sound. Because English's vowel
inventory is larger, multiple English vowels must collapse onto a single
Filipino letter: English \ipa{\ae} (``cat'') and \ipa{A:}
(``father'') are both written \gr{a}. This many-to-one relation creates an
ambiguity that English spelling alone cannot resolve, since the same letter
\gr{a} corresponds to \ipa{eI} in ``make'' but \ipa{\ae} in
``cat.'' Compounding this, English is stress-timed, so unstressed vowels reduce
to schwa \ipa{@}, a distinction that orthography does not encode at all.
Resolving which English vowel sound underlies a given letter therefore requires
information beyond spelling: the pronunciation itself.

The second property is one of spelling convention. The vowels (\gr{e},
\gr{i}, \gr{y}) and (\gr{o}, \gr{u}) are widely treated as
interchangeable in many lexical contexts, a convention rooted in the historical
allophonic status of these pairs in Filipino~\cite{schachter1972}. Native
speakers often accept either variant without perceiving an error, and this
latitude is especially common in borrowed vocabulary, where prescriptive
sources do not always fix a unique choice.

\subsection{Consonant Letter Correspondences and Context-Sensitive Rules}

For the consonant letters native to the Filipino alphabet~\cite{llamzon1966},
English--Filipino spelling correspondences are largely predictable one-to-one
mappings. However, English loanwords frequently contain sounds with no native
Filipino letter, notably \ipa{S}, \ipa{tS}, and \ipa{dZ}.
Prescriptive guidelines~\cite{almario2014,llamzon1966} supply standard written
forms: \ipa{tS} is written \gr{ts}, \ipa{dZ} as \gr{dy}, and
\ipa{S} as \gr{s} or \gr{sy}. Because such correspondences are fixed,
they can be applied as invariant substitutions.

However, not all consonant graphemes are invariant. Some are nativized
depending on adjacent graphemes.

\subsection{Spelling Repairs and Output Normalization}

Filipino spelling strongly favours syllables of the form consonant-vowel or
consonant-vowel-consonant~\cite{llamzon1966,magpale2024}. When a direct mapping
would yield a consonant cluster that violates this template, prescriptive
convention inserts a vowel letter to break the illegal
sequence~\cite{almario2014}. Enforcing this constraint requires a repair pass
over the mapped output.

